\documentclass{report}
\usepackage{fontspec}
\usepackage{xspace}
\usepackage{listings}
\usepackage{xcolor}
\usepackage{booktabs}
\usepackage{tabularx}
\usepackage[unicode,pdfencoding=auto]{hyperref}
\hypersetup{
    colorlinks=true,
    linkcolor=blue,
    urlcolor=cyan,
    pdfauthor={The Quarb Project},
    pdftitle={Quarb Implementation Guide},
}
\setmainfont{IBM Plex Serif}
\setmonofont{IBM Plex Mono}[Scale=MatchLowercase]

\definecolor{codegray}{rgb}{0.5,0.5,0.5}
\definecolor{backcolour}{rgb}{0.96,0.96,0.94}

\lstdefinelanguage{rust}{
  keywords={break, continue, else, for, if, in, loop, match, return, while,
    as, const, let, move, mut, ref, static, unsafe, use, where, crate, mod,
    pub, super, trait, impl, dyn, enum, struct, type, union, fn, extern, Self},
  keywordstyle=\color{blue}\bfseries,
  sensitive=true,
  comment=[l]{//},
  morecomment=[s]{/*}{*/},
  commentstyle=\color{codegray}\itshape,
  stringstyle=\color{orange},
  morestring=[b]",
}
\lstdefinestyle{mystyle}{
    backgroundcolor=\color{backcolour},
    basicstyle=\ttfamily\footnotesize,
    breaklines=true,
    captionpos=b,
    keepspaces=true,
    showstringspaces=false,
    tabsize=2,
    frame=none,
}
\lstset{style=mystyle}

\newcommand{\live}{\textbf{\textcolor{blue!60!black}{[live]}}\xspace}
\newcommand{\planned}{\textbf{\textcolor{gray}{[planned]}}\xspace}

\title{Quarb Implementation Guide}
\author{The Quarb Project}
\date{Version: Draft 0.1 (walking skeleton)}

\begin{document}

\maketitle
\tableofcontents

\chapter{About This Guide}

This guide describes the \textbf{Quarb engine}: the Rust
implementation of the Quarb query language, its \texttt{qua}
command-line tool, and the adapter architecture that lets one engine
query many data sources. The language itself\,---\,its data model,
execution model, and syntax\,---\,is defined in the \emph{Quarb
Language Specification} (the \texttt{lang} sibling project); this
guide covers only the implementation.

\paragraph{The guide tracks the code.}
The engine is being built as a series of thin, working slices rather
than from a finished design. This document tracks that process:
constructs the engine \emph{already implements} are marked \live and
documented against the real code; constructs that are specified but
\emph{not yet built} are marked \planned and sketch the intended
surface. As each slice lands, its \planned entries become \live. The
current milestone is a \emph{walking skeleton}: rooted tree
navigation over a filesystem.

\paragraph{Central concern: the adapter surface.}
The one design question this guide most needs to get right is the
\emph{adapter surface}\,---\,the trait a data source implements to be
queried. That surface is deliberately being grown against a concrete
first adapter (the filesystem) so that its shape is driven by real
navigation problems rather than speculation. Chapter~\ref{ch:adapter}
is its home.


\chapter{Architecture}

The engine follows a conventional compiler-style pipeline. A query
string is lexed and parsed into a small AST, which the executor
evaluates by driving an \emph{adapter} over the target data:

\begin{lstlisting}
[query string]
      |  lexer          quarb::lexer
      v
[tokens]
      |  parser         quarb::parser
      v
[query AST]             quarb::ast
      |  executor       quarb::exec
      v
[node set]  <-- drives --  [AstAdapter]  --> [data source]
\end{lstlisting}

\noindent
The executor never touches the data source directly. It knows only
the \texttt{AstAdapter} trait (Chapter~\ref{ch:adapter}), so the same
query runs unchanged over any adapter.

\section{Workspace Layout}

The engine is a Cargo workspace mirroring the intended separation of
concerns\,---\,the core engine, one crate per adapter, and the CLI:

\begin{tabularx}{\textwidth}{l X}
\toprule
\textbf{Crate} & \textbf{Responsibility} \\
\midrule
\texttt{quarb} & The core engine: \texttt{lexer}, \texttt{parser},
  \texttt{ast}, \texttt{exec}, and the \texttt{adapter} module
  defining \texttt{AstAdapter}. No I/O. \\
\texttt{quarb-fs} & The filesystem adapter. Depends on \texttt{quarb}. \\
\texttt{quarb-imap} & Remote IMAP mailboxes, queried in place:
  the maildir mapping over the wire (folders, messages by UID,
  headers as properties, \texttt{::;epoch}), transported by
  \texttt{curl} over subprocess (it speaks IMAP natively) — zero
  new dependencies. Lazy per message: one LIST, one UID SEARCH
  per folder, two fetches per touched message. The thread fabric
  is deliberately absent (indexing it means fetching every header
  — that is a sync; use mbsync + \texttt{quarb-maildir}). Target:
  \texttt{imap[s]://HOST[:PORT][/FOLDER]}; credentials via
  \texttt{QUARB\_IMAP\_USER}/\texttt{PASS} or
  \texttt{\textasciitilde/.netrc}. Read-only. Tested against a
  bundled deterministic mini IMAP server. \\
\texttt{quarb-json} & The JSON adapter (keys and indices as names,
  value types as traits). Depends on \texttt{quarb} and
  \texttt{serde\_json}. \\
\texttt{quarb-gsheet} & Google Sheets on the
  \texttt{quarb-xlsx} mapping (sheets as tables, first-row
  headers, typed values via \texttt{UNFORMATTED\_VALUE}),
  fetched from the Sheets API. Auth honesty: the API rejects
  gcloud's \texttt{cloud-platform} tokens, so the friction-free
  path is an API key (\texttt{?key=} /
  \texttt{QUARB\_GSHEET\_KEY}; enough for link-readable
  sheets), with a properly-scoped bearer for private ones.
  Target: \texttt{gsheet://SPREADSHEET\_ID}. \\
\texttt{quarb-html} & The HTML adapter (tags as names, block/inline
  as traits, attributes as properties, anchor fragments resolve).
  Depends on \texttt{quarb} and \texttt{scraper}. \\
\texttt{quarb-xml} & The XML adapter (qualified tags as names,
  attributes as properties, \texttt{id}/\texttt{xml:id} references
  resolve; strict well-formedness, no traits). Depends on
  \texttt{quarb} and \texttt{quick-xml}. \\
\texttt{quarb-relational} & The shared relational model: tables
  as root children, rows named by primary key, columns as
  properties, declared foreign keys driving
  \texttt{\textasciitilde>} resolution (chained hops instead of
  JOINs), \texttt{->}/\texttt{<-} crosslinks labeled by column,
  \texttt{<\textasciitilde} reverse resolution, and the
  table-naming hint for undeclared references. Loading is catalog-
  eager, rows-lazy: a table materializes on first touch through
  the driver-supplied fetcher, once, and untouched tables are
  never read (\texttt{::;loaded} reports the state). Engine
  adapters are thin catalog drivers over it. Depends on
  \texttt{quarb}. \\
\texttt{quarb-sqlite} & The SQLite catalog driver: PRAGMA
  introspection, storage-class type mapping. Depends on
  \texttt{quarb-relational} and \texttt{rusqlite}. \\
\texttt{quarb-objstore} & GCS and S3 buckets as lazy
  directory trees (\texttt{gs://} / \texttt{s3://}, delimiter
  listings, paginated; object content fetched on first read).
  Composed by default in \texttt{qua}: a bucket of JSON / CSV /
  source files grafts, so \texttt{/data/store.json/books/*}
  crosses from cloud storage into the document mid-path —
  grafting is this adapter's point. GCS: anonymous for public
  buckets, gcloud-token auth otherwise. S3: anonymous only in
  v1 (SigV4 signing is a recorded extension — private buckets
  refuse honestly). \\
\texttt{quarb-postgres} & The PostgreSQL catalog driver:
  \texttt{information\_schema} introspection over
  \texttt{tokio-postgres} (internal current-thread runtime,
  synchronous surface); native bool/int/float/text, everything
  else \texttt{::text}-cast into the numeric-reading regime.
  Depends on \texttt{quarb-relational}. \\
\texttt{quarb-mysql} & The MySQL/MariaDB catalog driver:
  \texttt{information\_schema} introspection over \texttt{sqlx}
  (internal current-thread runtime, synchronous surface;
  \texttt{key\_column\_usage} carries the foreign keys
  directly); integer/float/character families native, everything
  else \texttt{CAST(... AS CHAR)} into the numeric-reading
  regime. Depends on \texttt{quarb-relational}. \\
\texttt{quarb-csv} & The CSV/TSV adapter (one \texttt{row} node per
  record, columns as properties, empty cells as missing values that
  numeric aggregates skip; \texttt{::;columns} / \texttt{::;n-rows}
  metadata). Depends on \texttt{quarb} and \texttt{csv}. \\
\texttt{quarb-xpath} & The XPath~1.0 bridge, both directions:
  \texttt{translate} imports an XPath expression to Quarb, and
  \texttt{export} walks a Quarb query's reflection arbor (the
  locked vocabulary as the tooling surface) back to XPath —
  honest refusal and divergence notes both ways. Binaries
  \texttt{xpath2quarb} / \texttt{quarb2xpath}; the export side is
  differentially tested against \texttt{xmllint}. \\
\texttt{quarb-jq} & The jq bridge, both directions:
  \texttt{translate} imports a jq filter (value-projecting; a node
  mode pairs with \texttt{jq 'path(...)'}), and \texttt{export}
  walks the reflection arbor back to a jq filter. Binaries
  \texttt{jq2quarb} / \texttt{quarb2jq}; the export side is
  differentially tested against the \texttt{jq} binary. \\
\texttt{quarb-bigquery} & The Google BigQuery driver: the
  REST API directly (\texttt{jobs.query} over a synchronous HTTP
  client — no cloud SDK, no async runtime), bearer token from
  \texttt{QUARB\_BQ\_TOKEN} or \texttt{gcloud auth
  print-access-token}, catalog from
  \texttt{INFORMATION\_SCHEMA} including the \emph{unenforced}
  key constraints (declare them and the reference machinery
  works). The execution ladder is a billing instrument here:
  BigQuery meters bytes scanned, so lazy loading, pushdown, and
  \texttt{--save} directly reduce cost. Target syntax:
  \texttt{bigquery://PROJECT/DATASET[?account=EMAIL]}. \\
\texttt{quarb-duckdb} & DuckDB, the fifth catalog driver —
  and the columnar side door: a \texttt{CREATE VIEW ... AS
  SELECT * FROM read\_parquet(...)} makes Parquet first-class.
  Keys from \texttt{duckdb\_constraints()} (the
  \texttt{information\_schema} FK view reports the referencing
  side only), rows lazy, \texttt{raw\_query} for the ladder,
  read-only connection. Bundled build; \texttt{.duckdb} /
  \texttt{.ddb} by extension. \\
\texttt{quarb-xlsx} & Spreadsheets (\texttt{.xlsx} /
  \texttt{.xls} / \texttt{.ods}, via calamine) on the
  relational model: sheets as tables, first-row headers as
  columns (positional \texttt{cN} for blank/duplicate), rows
  named by sheet row number, numeric cells typed. No keys —
  sheets declare none — so resolution needs hints. Dispatched by
  extension \emph{before} the archive check (workbooks are
  zips; the sheets are the point). \\
\texttt{quarb-firebase} & The Firebase Realtime Database
  driver: the RTDB \emph{is} a JSON tree, so the mapping is
  \texttt{quarb-json}'s, fetched lazily one node per shallow
  GET and cached. Properties are direct fetches
  (\texttt{::score} is one tiny request); \texttt{\textasciitilde>}
  resolves by hint naming a root-relative container
  (\texttt{::parent\textasciitilde>item}). Containers that
  refuse enumeration (permission-scoped or unbounded — the
  public Hacker News database's \texttt{/v0/item} holds tens of
  millions of keys) are \emph{opaque}: reachable by name, empty
  to wildcards — served by the engine's new
  \texttt{children\_named} fast path. References are \emph{declared
  client-side}: a refs document
  (\texttt{\{"refs": \{"parent": "item", "kids/*": "item"\}\}},
  supplied via \texttt{qua --refs}) is the portable mini-catalog
  a schemaless substrate cannot hold — with it, bare
  \texttt{\textasciitilde>} resolves and \texttt{->} crosslinks
  enumerate (array fields one probed edge per element); an inline
  hint overrides. Reverse resolution stays empty (it would scan
  what opaque containers refuse; server-side \texttt{.indexOn}
  queries are the recorded v2 path). Target:
  \texttt{firebase://HOST/PATH[?auth=...]}. Read-only. \\
\texttt{quarb-code} & Source code as its syntax tree
  (tree-sitter): node kinds as edge names
  (\texttt{//function\_item}), tree-sitter \emph{fields} as
  properties (\texttt{::name} on a function is its name),
  source text as the value, \texttt{::;start-line} /
  \texttt{::;kind} / \texttt{::;field} metadata; named nodes
  only. Rust / Python / JavaScript by extension, compile-time
  grammar set. Grafts under composition — \texttt{--descend}
  queries every function in a directory tree at once. \\
\texttt{quarb-compose} & Adapter \emph{composition}: wrap any
  adapter and a leaf whose name or content parses (json / xml /
  html / csv, else a sniff) grafts its parsed arbor as children —
  \texttt{/data/store.json/books/*} crosses from filesystem into
  JSON mid-path. Grafts are lazy (read + parse on first entry),
  the inner root is identified with the outer leaf (no phantom
  node), and locators mark the boundary jar-style
  (\texttt{store.json!/books/1}). \texttt{qua}: default for
  archives, \texttt{--descend} for directories. \\
\texttt{quarb-archive} & Zip and tar (\texttt{.tar[.gz]})
  archives as path trees: entry content as the node value,
  \texttt{::;size} / \texttt{::;compressed}; zip contents lazy
  per entry, tar in one pass. With composition, \texttt{.docx} /
  \texttt{.epub} / \texttt{.jar} internals are directly
  queryable. \\
\texttt{quarb-maildir} & Maildir (\texttt{cur/}+\texttt{new/})
  and mbox mailboxes: headers as properties (lowercased), the
  body as the value, \texttt{::;epoch} parsing the Date header
  for range predicates. The thread graph is reference fabric:
  \texttt{::in-reply-to\textasciitilde>} walks to the parent
  message, \texttt{<-in-reply-to} finds the replies, via a
  Message-ID index. Never touches flags; read-only. \\
\texttt{quarb-firestore} & Google Firestore (native mode):
  collections and documents alternate as the arbor's levels; map
  and array fields descend as children, subcollections beside
  them. \texttt{referenceValue} fields are \emph{native}
  references — the first remote adapter where
  \texttt{\textasciitilde>} and \texttt{->} need neither a
  schema nor a refs file. Lazy per level, REST direct, token via
  \texttt{QUARB\_GCP\_TOKEN} / gcloud. Target:
  \texttt{firestore://PROJECT/DATABASE[?account=...]}. \\
\texttt{quarb-datastore} & Google Cloud Datastore (Firestore in
  Datastore mode): kinds and entities, the most relational of the
  document stores; \texttt{keyValue} properties are native
  references (\texttt{\textasciitilde>}/\texttt{->} schema-free).
  Kind list via a \texttt{\_\_kind\_\_} query, entities via
  cursor-paginated \texttt{runQuery}, lazy per kind. Target:
  \texttt{datastore://PROJECT[/DATABASE][?account=...]}. Both
  drivers send the \texttt{x-goog-request-params} routing header
  named databases require. \\
\texttt{quarb-git} & The git repository driver: the object
  DAG as the arbor. \texttt{/branches}, \texttt{/tags},
  \texttt{/HEAD} (refs alias their commits — same properties,
  same children), \texttt{/commits} (rev-syntax aliasing:
  \texttt{/commits/'HEAD\textasciitilde2'} navigates without
  enumeration; \texttt{/commits/*} batches one
  \texttt{rev-list --all}). A commit's children are its
  \emph{tree} — time-travel file access, blob content as the
  node value. References are inherent:
  \texttt{::parent\textasciitilde>},
  \texttt{->parent} (merges fan out), \texttt{<-parent}
  (children, from the enumeration). The diff surface: tree
  entries carry the \texttt{<changed>} trait when their path is
  in the commit's diff vs its first parent (trees by prefix), so
  \texttt{git log -- path} is
  \texttt{/commits/*[/path<changed>]}; commits answer
  \texttt{::changed} (the touched paths, deletions included)
  and \texttt{::;n-changed}. Plumbing over subprocess —
  no libgit2, zero new dependencies. Target: \texttt{git:PATH}.
  Read-only. \\
\texttt{quarb-serve} & The \emph{serve protocol} — the
  extension surface for tools the engine never links: a tool
  implements \texttt{AstAdapter} over its own model and hands it
  to \texttt{serve()} (three lines; non-Rust tools implement the
  protocol directly — nine operations over line-oriented JSON),
  and \texttt{qua serve:COMMAND} spawns it and mounts it beside
  every other source. LSP's shape: the protocol is the plugin
  system. Two wires, one protocol: the handshake is JSON (the
  universal floor) and advertises formats; every first-party
  server offers \emph{kaiv} (\texttt{.daiv} frames, natively
  typed leaves under \texttt{/@serve}) and the client upgrades
  to it — the house format is the default wire, JSON the
  foreign-tool floor. Read-only by construction (no mutating
  operation exists). First external user: \texttt{cuj quarb-serve}
  (the cuj vault, mounted without qua knowing cuj). Reference
  server: \texttt{quarb-serve-json}; identity-tested against the
  direct adapter. \\
\texttt{quarb-sql} & The SQL bridge, both directions:
  \texttt{translate} imports a \texttt{SELECT} statement (WHERE,
  one JOIN via correlation and the witness, GROUP BY/HAVING,
  ORDER/LIMIT/DISTINCT, \texttt{LIKE} as case-insensitive
  regexes), and \texttt{export} walks a Quarb query's reflection
  arbor back to \texttt{SELECT} — witness joins become
  \texttt{JOIN ... ON}, grouped reductions become \texttt{GROUP
  BY}/\texttt{HAVING}; \texttt{\textasciitilde>} chains are
  refused by design (foreign keys live in the schema, not the
  query — spell \texttt{<=>} to export). Both directions
  differentially tested against SQLite; binaries
  \texttt{sql2quarb} / \texttt{quarb2sql}. \\
\texttt{qm} & The mailbox interface for the terminal — the
  first \emph{library-path} client beside \texttt{qua}: it opens
  the Maildir adapter once per invocation, addresses messages as
  node handles, and phrases its subcommands (\texttt{ls},
  \texttt{read}, \texttt{thread}, \texttt{from}, \texttt{grep},
  \texttt{count}) as queries where a query is the natural shape,
  with \texttt{qm query} as the raw escape hatch. Actions are the
  tool's, not the engine's: \texttt{qm seen} performs the Maildir
  flag rename itself on the file handle the adapter exposes
  (\texttt{path\_of}). \\
\texttt{qua} & The command-line tool. Picks the adapter from its
  argument (a directory is filesystem; \texttt{.csv}/\texttt{.tsv}
  is a table; a file or piped stdin is XML by extension or a leading
  \texttt{<?xml}, HTML by extension or a leading \texttt{<}, else
  JSON) and prints results. \texttt{--xpath} / \texttt{--jq} / \texttt{--sql}
  translate the query from XPath, jq, or SQL first; database
  inputs get \emph{pushdown} — a query whose constructs are all
  in the verified-safe set (filtered selections, whole-table
  aggregates, witness joins) compiles to SQL and runs on the
  server, with rows ordered by the catalog key so output is
  identical to the scan path (\texttt{--no-pushdown} opts out;
  any refusal or error falls back silently, and
  \texttt{--explain} prints the decision — the pushed SQL, or the
  construct that kept the query on the scan path). When full
  pushdown refuses, \emph{partial} pushdown pushes the leading
  row predicates as the table's fetch filter and runs the original
  query over the reduced set (gated: single reach, no
  crosslink/resolution axes, no metadata — anything that could
  observe the filtering). \texttt{--save FILE} materializes a
  result — \texttt{.db}/\texttt{.sqlite} as a table
  (\texttt{--as} names it), \texttt{.json} as a JSON array —
  both first-class inputs for later queries (the reference syntax
  is the existing file/mount mechanism); \texttt{--daiv} emits
  canonical kaiv (typed leaves with per-value provenance; needs the
  sibling \texttt{kaiv-rs} checkout). \\
\bottomrule
\end{tabularx}

\medskip

Adapters are separate crates by design: \texttt{quarb-fs},
\texttt{quarb-json}, \texttt{quarb-html}, and \texttt{quarb-xml}
share nothing but the \texttt{AstAdapter} trait, and a consumer
pulls in only the adapters it needs. Each new adapter has confirmed
the surface without changing it\,---\,the filesystem, a JSON tree,
an HTML DOM, and an XML document are genuinely the same arbor to
the engine.

\section{The Execution Ladder}\label{sec:ladder}

Database adapters evaluate on a three-rung ladder; every rung is
\emph{output-identical} to the one below it (differential tests
hold this), so the rungs trade only latency and transfer:

\begin{enumerate}
  \item \textbf{Full pushdown} (\texttt{quarb\_sql::pushdown},
    the exporter's strict mode): the whole query compiles to SQL
    when every construct is in the verified-safe set —
    WHERE-filtered selections, whole-table aggregates, witness
    joins. Strictness is the point: anything whose SQL semantics
    are not provably identical (LIKE case folding, truthiness,
    GROUP BY/DISTINCT order, collations, unordered LIMIT) refuses.
    Row order is part of the contract — the plan names the result
    table, and each driver appends \texttt{ORDER BY} its key from
    the catalog. Measured: COUNT over 500k PostgreSQL rows,
    2.57\,s scanned vs.\ 0.018\,s pushed.
  \item \textbf{Partial pushdown}
    (\texttt{quarb\_sql::partial\_pushdown}): the leading run of
    strictly-translatable row predicates becomes the table's fetch
    filter (each driver's lazy fetcher accepts a per-table WHERE),
    and the \emph{original} query runs unchanged over the reduced
    set — the engine re-applies the pushed predicates, a no-op on
    pre-filtered rows, so there is no query rewriting. Gates:
    leading run only (a positional predicate before an expression
    one sees unfiltered rows), a single \texttt{/table/*} branch
    with no correlations, the table reached exactly once, no
    crosslink/resolution axes, no \texttt{:::}/\texttt{::;}
    metadata (each would observe the filtering). Measured:
    unpushable grouping over a 1\%-selective filter, 2.31\,s
    vs.\ 0.081\,s.
  \item \textbf{Lazy scan}: catalog-eager, rows-lazy per table
    (Section on \texttt{quarb-relational}); correct for
    everything, always the fallback.
\end{enumerate}

\noindent
\texttt{qua --explain} prints the decision chain; \texttt{qua
--no-pushdown} pins the scan rung; \texttt{qua --save FILE
[--as NAME]} materializes any result as a first-class input
(\texttt{.db} table / \texttt{.json} array) so expensive
reductions run once and iteration proceeds locally. For metered
engines (a future BigQuery adapter bills by bytes scanned), the
ladder's transfer reductions are cost reductions.

\section{Current Query Subset}

The engine implements essentially the whole language. A query is a
root-anchored sequence of hops, each an axis followed by a name:

\begin{tabularx}{\textwidth}{l X}
\toprule
\textbf{Form} & \textbf{Meaning} \\
\midrule
\texttt{/name}, \texttt{//name} & \live child / descendant named \texttt{name} \\
\texttt{\textbackslash name}, \texttt{\textbackslash\textbackslash name}
  & \live parent / ancestor named \texttt{name} \\
\texttt{>name}, \texttt{<name} & \live next / previous sibling \\
\texttt{>>name}, \texttt{<<name}, \texttt{>>?}, \texttt{>>!},
\texttt{<<?}, \texttt{<<!}
  & \live sibling reach: all following / preceding siblings at any
    distance, with nearest / farthest quantifiers (matcher before
    reach, mirroring \texttt{//?}); covers CSS \texttt{\textasciitilde}
    and XPath following/preceding-sibling \\
\texttt{//?}, \texttt{//!}, \texttt{\textbackslash\textbackslash?}, \texttt{\textbackslash\textbackslash!}
  & \live proximal (nearest) / distal (farthest) reach \\
\texttt{\textasciicircum}, \texttt{name\$} & \live root anchor / leaf anchor \\
\texttt{*}, \texttt{*.rs}, \texttt{\textasciitilde(re)}
  & \live wildcard, glob, and regex name matching \\
\texttt{<code>}, \texttt{<a|b>}, \texttt{<a><b>}
  & \live trait filters (alternation, conjunction) \\
\texttt{::}, \texttt{:::key}, \texttt{::;key}
  & \live property, core-metadata, and adapter-metadata projection \\
\texttt{[::;size > 1K]}, \texttt{[/child]}, \texttt{[n]},
\texttt{[-1]}, \texttt{[2..5]}
  & \live predicates: comparisons (\texttt{=} \texttt{!=} \texttt{<}
    \texttt{<=} \texttt{>} \texttt{>=}), substring \texttt{*=}, regex
    \texttt{=\textasciitilde}/\texttt{!\textasciitilde}, structural
    conditions, \texttt{and}/\texttt{or}/\texttt{!}/\texttt{not},
    positional index and range (negative counts from the end) \\
\texttt{::price * ::qty}, \texttt{7 div 2}, \texttt{idiv}, \texttt{mod}
  & \live value expressions: arithmetic over numeric readings, in
    predicates, pipeline stages (\texttt{| ::a * ::b}), and pushes
    (\texttt{| .total(::a * ::b)} — a computed column); null
    propagation, exact integers, overflow to float \\
\texttt{//a || //b}, \texttt{| count}, \texttt{| sort | join(",")}
  & \live union, and pipelines over a standard-library function set \\
\texttt{->target}, \texttt{<-target}, \texttt{->*}
  & \live crosslink navigation (filesystem: symlinks) \\
\texttt{(/ul/li)+}, \texttt{(/p|/div)}, \texttt{(->mgr)\{1,3\}!},
\texttt{/\{2\}}
  & \live path patterns: groups, subpath alternation (strict and
    tolerated forms), quantifiers (\texttt{+} \texttt{*}
    \texttt{\{n\}} \texttt{\{m,n\}} \texttt{\{m,\}}) with
    proximal/distal suffixes, the pattern dot wildcard, and the
    bare-operator sugar. Simple-path expansion (no node revisited
    within one path, the anchor included), capped by the adapter's
    quantifier bound (default 32; \texttt{--quantifier-bound}
    overrides) \\
\texttt{| upper}, \texttt{@| count}, \texttt{| .n}, \texttt{| \$.n}
  & \live per-capsa transform, whole-context aggregation, register
    push and recall \\
\texttt{@| sort\_by(::age)}, \texttt{@| top(3, ::fare)},
\texttt{@| [..n]}, \texttt{| [::age > 30]}
  & \live keyed aggregates (\texttt{sort\_by}, \texttt{unique\_by},
    \texttt{min\_by}/\texttt{max\_by}, \texttt{top}/\texttt{bottom};
    capsa-preserving, missing keys do not compete), whole-context
    positional selection, and per-capsa filters \\
\texttt{@| group(::k) | top(2, ::v) @| ungroup}, \texttt{\$ord}
  & \live value-keyed grouping (list topics + member capsae, keys as
    named regs), per-group keyed stages on the plain pipe,
    \texttt{ungroup} flattening, and the \texttt{\$ordinal} operand
    (position in the current context, scope-threaded into filters) \\
\texttt{| \%.}, \texttt{qua q.quarb}
  & \live the register record view (named regs as a record: the
    row-wise table view, subsuming the planned \texttt{@.|}) and
    reflection: \texttt{.quarb} files open as the query's own arbor
    (\texttt{quarb::reflect}; vocabulary locked v1, enforced by the
    conformance corpus in \texttt{tests/vocabulary.rs}, version on
    the root as \texttt{::;vocabulary}) \\
\texttt{def \&f(\$x): ...;}, \texttt{\&f(::age)},
\texttt{--defs}, \texttt{--expand}
  & \live fragments: parse-time expansion with call-by-name
    \texttt{\$param} forms, defs files, and the unparser
    (\texttt{quarb::unparse}) rendering any parsed query back to
    canonical text (\texttt{--expand} = macroexpand) \\
\texttt{| [::x =\textasciitilde{} /(\textbackslash{}w+)/] | rec("k", \$1)},
\texttt{qua -i}
  & \live match captures (filter-stage \texttt{=\textasciitilde}
    binds groups onto the capsa) and the interactive session
    (living-query lines with rollback, session defs,
    \texttt{:show}/\texttt{:undo}) \\
\texttt{@| window(-2..0) | mean}, \texttt{@| window(..0) | sum},
\texttt{@| shift(1, ::k)}
  & \live window functions: offset-range spans over context
    neighbors as members (partial edges; \texttt{window(n)} =
    trailing count; open start = expanding), \texttt{shift} lag
    with own registers kept, and optional partition keys
    (per-group windows in original row order) \\
\texttt{"\$\{::name\} (\$\{::age\})"},
\texttt{macro \&m(\$p, @rest): body;}
  & \live string interpolation (double-quoted \texttt{\$\{expr\}}
    holes, evaluated in the current scope; null splices empty) and
    procedural macros: the body runs at expansion time against the
    invocation's expansion arbor (argument forms as reflected
    subtrees, \texttt{::form} = each form's text), literal args bind
    by value, \texttt{@rest} collects the remainder, and the text
    results — joined, reparsed against strictly-earlier fragments —
    become the expansion (\texttt{--expand} = macroexpand) \\
\texttt{macro \&pivot!(\$col): /data/row | \$col @| unique | ...}
  & \live data-aware macros: \texttt{!} on the definition mounts the
    queried dataset as \texttt{/data} in the expansion arbor, and
    every invocation spells the \texttt{!} (enforced both ways);
    outside a hole a parameter is the argument form (call-by-name),
    inside a hole its text — so a pivot reads distinct data values
    at expansion and generates a column per value;
    \texttt{--expand} with an input expands against it \\
\texttt{| .(//*.rs @| count)}, \texttt{| .t(\^{}/row @| count)}
  & \live subcontexts (grouped aggregation); a \texttt{\^{}}-anchored
    branch navigates from the arbor root, so a subcontext can
    consult the whole table (share-of-total columns) \\
\texttt{.c(\^{}/row[::k = \$\$.k] @| count)},
\texttt{| rec(\$*1/name::, /amt::)}
  & \live correlated subqueries: \texttt{\$\$} steps a capsa-scope
    operand out to the invoking capsa (\texttt{\$\$.name},
    \texttt{\$\$\_}, \texttt{\$\$ord}; stacking
    \texttt{\$\$\$} steps further) — wide crosstabs, groupwise
    transforms; and the correlation keeps its \emph{witness}, so
    \texttt{\$*k} in pipeline stages projects the joined side into
    output (\texttt{SELECT a.x, b.y}) \\
\texttt{E1 <=> E2[/k:: = \$*1/k:: and …]}
  & \live correlation (context join): one tuple binds the whole
    predicate, and \texttt{\$*N} may navigate from its context node
    (\texttt{\$*1/field::}), so joins work where keys are children,
    e.g. JSON \\
\texttt{::ref\textasciitilde>}, \texttt{::'\$ref'\textasciitilde>/city}
  & \live cross-reference resolution (JSON: follow a \texttt{\$ref}
    JSON Pointer) \\
\texttt{::ref<\textasciitilde}
  & \live reverse resolution: every node whose \texttt{ref} resolves
    to this one (naive whole-arbor scan) \\
\bottomrule
\end{tabularx}

\medskip

The pipeline runs on \emph{capsae}\,---\,the register system's unit,
a node paired with a register (its breadcrumbs) and a current topic.
A per-capsa \texttt{|} transforms each topic; \texttt{@|} reduces the
whole context; \texttt{| .name} pushes the topic onto the register and
\texttt{| \$.name} recalls it; a subcontext \texttt{| .(expr)}
evaluates \texttt{expr} from each node, reduces it to a scalar, and
pushes it\,---\,grouped aggregation without leaving the query.

\medskip

Any other operator is lexed and rejected with a
\texttt{not yet supported} error naming the offending construct,
rather than silently mis-parsed. Everything in
Section~\ref{sec:planned} falls into this category today.


\chapter{The Adapter Surface}\label{ch:adapter}

An \emph{adapter} maps a data source onto the Quarb arbor model\,---\,a
tree backbone whose edges carry \emph{names}, enriched with
non-hierarchical crosslinks. The engine drives all navigation through
the \texttt{AstAdapter} trait, so this trait \emph{is} the contract
between the language and the world.

\section{The Live Trait}

The methods the skeleton actually uses. Signatures are the real ones
from \texttt{quarb/src/adapter.rs}:

\begin{lstlisting}[language=rust]
/// An opaque handle to a node, minted and interpreted by the
/// adapter that produced it.
pub struct NodeId(pub u64);

pub trait AstAdapter {
    /// The root node -- the initial navigation context.
    fn root(&self) -> NodeId;

    /// The tree children of `node`, in document order.
    fn children(&self, node: NodeId) -> Vec<NodeId>;

    /// The name of `node` -- the label of its incoming tree edge.
    /// `None` when the adapter leaves a node unnamed (e.g. the
    /// filesystem root).
    fn name(&self, node: NodeId) -> Option<String>;

    /// The parent of `node`, or `None` for the root.
    fn parent(&self, _node: NodeId) -> Option<NodeId> { None }

    // --- projection surface (defaulted; implement what fits) ---

    /// Adapter-defined classifications, for `<trait>` filters.
    fn traits(&self, _node: NodeId) -> Vec<String> { Vec::new() }

    /// A named property `::prop`.
    fn property(&self, _node: NodeId, _name: &str) -> Option<Value> { None }

    /// The default projection, bare `::` (filesystem: file content).
    fn default_value(&self, _node: NodeId) -> Option<Value> { None }

    /// Adapter metadata `::;key` (filesystem: size, modified, …).
    fn metadata(&self, _node: NodeId, _key: &str) -> Option<Value> { None }

    /// Outgoing / incoming crosslinks as (label, node) pairs, for
    /// `->` / `<-` (filesystem: symlinks).
    fn links(&self, _node: NodeId) -> Vec<(String, NodeId)> { Vec::new() }
    fn backlinks(&self, _node: NodeId) -> Vec<(String, NodeId)> { Vec::new() }
}
\end{lstlisting}

\noindent
The navigation methods are required; the projection methods carry
defaults, so an adapter implements only what its domain offers. Core
metadata (\texttt{:::key}\,---\,index, depth, counts, \texttt{is-leaf},
\dots) is computed by the engine from the navigation surface, so no
adapter implements it.

Three observations that the filesystem work has already forced:

\begin{itemize}
\item \textbf{Names live on edges, not nodes.} \texttt{name} returns
  the label of a node's \emph{incoming} edge\,---\,its role in its
  parent\,---\,exactly as the specification's edge-name model
  requires. For the filesystem this is simply the filename.

\item \textbf{\texttt{NodeId} is opaque and adapter-private.} The
  engine never inspects the \texttt{u64}; it only passes ids back to
  the adapter. Each adapter chooses its own identity scheme (the
  filesystem adapter uses an interning arena index).

\item \textbf{The root may be unnamed.} \texttt{name} returns
  \texttt{Option} precisely so an adapter can leave the root without a
  name, as the filesystem root \texttt{/} does.
\end{itemize}

\section{Worked Example: the Filesystem Adapter}

\texttt{quarb-fs} implements the live trait over a directory tree. It
materializes the tree once, at construction, into arrays indexed by
\texttt{NodeId} (path, children, parent); the trait methods are then
trivial lookups:

\begin{lstlisting}[language=rust]
impl AstAdapter for FsAdapter {
    fn root(&self) -> NodeId { self.root }
    fn children(&self, node: NodeId) -> Vec<NodeId> { self.children[node].clone() }
    fn name(&self, node: NodeId) -> Option<String> {
        if node == self.root { return None; }   // "/" is unnamed
        self.paths[node].file_name().map(|n| n.to_string_lossy().into())
    }
    fn parent(&self, node: NodeId) -> Option<NodeId> { self.parents[node] }
}
\end{lstlisting}

\paragraph{Ignore rules \live.}
The materializing walk uses \texttt{ignore::WalkBuilder}\,---\,the
same machinery ripgrep uses\,---\,so by default the adapter respects
\texttt{.gitignore}/\texttt{.ignore} files and skips hidden entries.
This came directly from running the skeleton: \texttt{qua '//*.rs'}
in the engine's own tree first returned a build artifact under
\texttt{target/}, because nothing yet excluded it. The adapter now
honors ignore rules by default, and \texttt{qua} exposes
\texttt{-{}-hidden} and \texttt{-{}-no-ignore} to opt back in. The
walk is currently eager (a targeted query still walks the whole
tree); a lazy traversal with cascading ignore rules is a planned
optimization. This is the pattern the guide follows: a concrete
problem, surfaced by running the code, drives the next slice.

\section{The Planned Surface}\label{sec:planned}

The rest of the arbor model, specified in the language but not yet
built. Each becomes a live trait method (or a related capability) as
its engine slice lands. The intended shapes:

\begin{tabularx}{\textwidth}{l l X}
\toprule
\textbf{Capability} & \textbf{Query} & \textbf{Adapter surface (sketch)} \\
\midrule
Pattern search & \texttt{=>}
  & \texttt{search(pattern) -> Vec<NodeId>}. \\
\bottomrule
\end{tabularx}

\medskip

These signatures are provisional. They are recorded here so the
direction is legible, not to freeze a contract\,---\,the filesystem
adapter (and the JSON adapter after it) will reshape them as they are
implemented.


\chapter{Building and Running}

\begin{lstlisting}[language=bash]
# from the engine workspace root
cargo build            # build the whole workspace
cargo test             # unit + integration + doc tests

# query a directory tree
cargo run -p qua -- '//*.rs' src/     # all Rust files under src/
cargo run -p qua -- '/README.md'      # a single top-level file

# query a document
cargo run -p qua -- '/users/0/name::' data.json
cargo run -p qua -- '//book[::pages > 200]::;id' catalog.xml
cargo run -p qua -- '//a::href' page.html
\end{lstlisting}

\noindent
\texttt{qua <QUERY> [PATH]} runs \texttt{QUERY} against the directory
\texttt{PATH} (default: the current directory) and prints each
matching path on its own line, so it composes in shell pipelines
alongside \texttt{sort}, \texttt{uniq}, and friends.


\chapter{Planned Integration Paths}

Beyond the direct \texttt{AstAdapter} implementation shown above,
these integration paths are described in the specification but
\planned\ (none is implemented yet):

\begin{itemize}
\item \textbf{Petgraph adapter.} A ready-made adapter over a
  \texttt{petgraph::Graph}, for Rust callers who transform their AST
  into a graph once and query it at speed. See
  \texttt{petgraph-adapter.md} in this project for the design note.
\item \textbf{Python bindings.} A \texttt{quarb} Python package
  exposing the engine and an \texttt{AstAdapter} base class.
\item \textbf{C FFI.} A C ABI for embedding the engine in other
  languages.
\end{itemize}

\noindent
Each will be added once the core language surface is substantial
enough to be worth binding.

\end{document}
